\documentclass[11pt]{article}
\usepackage[a4paper,margin=30mm]{geometry}
\usepackage{fontspec}
\setmainfont{DejaVu Serif}
\usepackage{amsmath,amssymb}
\usepackage{microtype}
\usepackage{enumitem}
\usepackage{hyperref}
\hypersetup{colorlinks=false,pdfborder={0 0 0}}
\newcommand{\leg}[2]{\left(\frac{#1}{#2}\right)}
\newcommand{\unitid}[1]{\hypertarget{#1}{}}
\setlength{\parindent}{1.2em}
\setlength{\parskip}{0.25em}
\begin{document}
\begin{center}
{\Large\bfseries RESEARCHES ON THE PRIME DIVISORS OF A CLASS\\
OF FORMULAE OF THE FOURTH DEGREE.}\\[1.2em]
{First memoir. Continuation: printed pages 75--80.}
\end{center}

\unitid{dirichlet.v1.p04b.en.par.001}It is very easy to make sure that the numbers $E$ and $F$ are relatively prime. Indeed, let $\delta$ be any simple factor of $E$ (which will necessarily be odd, $E$ itself being odd), and suppose that $\delta$ also divides $F$. Inspection of the preceding equations shows that, under this supposition, $\delta$ would be a common divisor of $t+\psi$ and $t-\psi$, and consequently would divide the number $t$, which is the half-sum of the preceding two. But on the other hand one sees that $\delta$, as a prime divisor of $E$, divides $m^2$, and consequently $m$ and $u$ ($u$ being $=mu'$). The numbers $t$ and $u$ would therefore, under this supposition, have the common factor $\delta$, which is absurd, since $p=t^2-bu^2$ is a prime number. It is therefore proved that $E$ and $F$ cannot have a common divisor; and since the product of these numbers is a square, each of them is likewise a square. Thus one has:
\unitid{dirichlet.v1.p04b.en.eq.001}
\[
\leg{E}{b}=1,\qquad \leg{F}{b}=1.
\]

\unitid{dirichlet.v1.p04b.en.par.002}As for the numbers $K$ and $L$, they are likewise such that one has:
\unitid{dirichlet.v1.p04b.en.eq.002}
\[
\leg{K}{b}=1,\qquad \leg{L}{b}=1.
\]

\unitid{dirichlet.v1.p04b.en.par.003}Indeed, one sees from the last equation (p. 74) that these numbers divide, the one and the other, the formula $\varphi'^2+bu'^2$ (where $\varphi'$ and $bu'$ are relatively prime); and it is known, by a known theorem that follows easily from the law of reciprocity (\emph{Théorie des Nombres}, no. 197), that every odd divisor $R$ of such a formula is such that $\leg{R}{b}=1$. Multiplying the first of the last relations by the third, and the second by the fourth, one obtains:
\unitid{dirichlet.v1.p04b.en.eq.003}
\[
\leg{EK}{b}=1,\qquad \leg{FL}{b}=1,
\]
or, which is the same thing:
\unitid{dirichlet.v1.p04b.en.eq.004}
\[
(\delta)\qquad \leg{t+\psi}{b}=1,\qquad \leg{t-\psi}{b}=1.
\]

\unitid{dirichlet.v1.p04b.en.par.004}Examining in a similar way the case in which $\varphi$ is divisible by $b$, one will find that, in this case, one of the numbers $t+\psi$ and $t-\psi$ is divisible by $b$, and that the other is, as in the case with which we have just been occupied, a quadratic residue with respect to $b$; so that, denoting by $t\pm\psi$ that one of these numbers which is divisible by $b$, one has:
\unitid{dirichlet.v1.p04b.en.eq.005}
\[
(\delta')\qquad t\pm\psi\equiv0\pmod b,\qquad
\leg{t\mp\psi}{b}=1.
\]

\unitid{dirichlet.v1.p04b.en.par.005}It follows from the theorem stated at the end of the preceding paragraph that, in order to decide whether $-b$ is or is not a biquadratic residue with respect to $p$, everything reduces to the question of knowing whether $t$ is or is not a quadratic residue with respect to $b$. We can now show, by means of the results just obtained, that this latter question can be decided without knowing $t$, that is, without its being necessary to put $p$ under the form $t^2-bu^2$. Suppose first that $\varphi$ is not divisible by $b$. The relations $(\delta)$, which hold in this case, can be presented in this way, observing that $b$ is a prime number $4n+3$:
\unitid{dirichlet.v1.p04b.en.eq.006}
\[
\leg{\psi+t}{b}=1,\qquad \leg{\psi-t}{b}=-1.
\]

\unitid{dirichlet.v1.p04b.en.par.006}One immediately draws from the equation $p=t^2-bu^2$: $t^2\equiv p\pmod b$, a result which shows that, on solving the congruence $\chi^2\equiv p\pmod b$, and denoting by $\chi$ one of its roots taken at random, one will have $t\equiv\chi$, or $t\equiv-\chi\pmod b$. In the first case, one has:
\unitid{dirichlet.v1.p04b.en.eq.007}
\[
\leg{t}{b}=\leg{\chi}{b},\qquad \leg{\chi+\psi}{b}=1,
\]
for it is evident that the expressions $\leg{t}{b}$ and $\leg{t+\psi}{b}$ do not change when one puts, in place of $t$, another number $\chi$ which differs from $t$ only by a multiple of $b$. Multiplying the last relations together, one obtains:
\unitid{dirichlet.v1.p04b.en.eq.008}
\[
\leg{t}{b}=\leg{\chi(\chi+\psi)}{b}.
\]

\unitid{dirichlet.v1.p04b.en.par.007}Consider now the other case, in which $t\equiv-\chi\pmod b$. One will then have:
\unitid{dirichlet.v1.p04b.en.eq.009}
\[
\leg{t}{b}=-\leg{\chi}{b},\qquad \leg{\psi+\chi}{b}=-1,
\]
the last of these relations resulting from the formula $\leg{\psi-t}{b}=-1$, if one puts, in place of $-t$, the number $\chi$, which differs from it only by a multiple of $b$. Multiplying the preceding relations together, one will find, as above:
\unitid{dirichlet.v1.p04b.en.eq.010}
\[
\leg{t}{b}=\leg{\chi(\chi+\psi)}{b}.
\]

\unitid{dirichlet.v1.p04b.en.par.008}We have therefore arrived at this remarkable result: if $\varphi$ is not divisible by $b$, one can decide very simply whether one has:
\unitid{dirichlet.v1.p04b.en.eq.011}
\[
\leg{t}{b}=1\quad\hbox{or}\quad \leg{t}{b}=-1.
\]

\unitid{dirichlet.v1.p04b.en.par.009}It will suffice to seek a number satisfying the congruence $\chi^2\equiv p\pmod b$; having found such a number $\chi$, one will always have:
\unitid{dirichlet.v1.p04b.en.eq.012}
\[
\leg{t}{b}=\leg{\chi(\chi+\psi)}{b}.
\]

\unitid{dirichlet.v1.p04b.en.par.010}Let us now come to the case in which $\varphi$ is divisible by $b$. The relations $(\delta')$ which then hold are:
\unitid{dirichlet.v1.p04b.en.eq.013}
\[
t\pm\psi\equiv0\pmod b,\qquad
\leg{t\mp\psi}{b}=1.
\]
In the last of these relations one may add to $t\mp\psi$ any multiple of $b$. If one adds to it the number $t\pm\psi$, which, by virtue of the first relation, is divisible by $b$, there will result $\leg{2t}{b}=1$, a result which entails this other one: $\leg{t}{b}=\leg{2}{b}$.

\unitid{dirichlet.v1.p04b.en.par.011}Combining this result with a known theorem, one will see that, in the case of $\varphi$ divisible by $b$, one has $\leg{t}{b}=1$ or $\leg{t}{b}=-1$, according as $t$ is of the form $8n+7$ or of this one: $8n+3$.

\unitid{dirichlet.v1.p04b.en.par.012}Comparing what has just been proved with the theorem established at the end of the preceding paragraph, one will see that it is possible to decide, independently of the knowledge of the number $t$, whether $-b$ is or is not a biquadratic residue with respect to $p$. Since the result thus reached no longer contains any trace of the number $t$, one is led to believe that it does not presuppose the possibility of the equation $t^2-bu^2=p$, from which the number $t$ takes its origin, and that it is generally true for all values of $b$ and $p$ such that $\leg{b}{p}=1$. This is indeed what takes place, as one can prove, by examining, instead of the equation $t^2-bu^2=p$, the more general equation $t^2\pm Mu^2=p$, where $M$ denotes the product of any number of distinct prime numbers, and then combining the results of this examination with the following theorem, of whose truth one cannot doubt, but whose rigorous demonstration nevertheless presents difficulties:

\unitid{dirichlet.v1.p04b.en.par.013}“$c$ denoting any prime number whatever and $p$ a prime number $4n+1$, such that $\leg{c}{p}=1$, one can always determine a number $\delta$, composed of simple factors all smaller than $c$, and such that the equation $t^2\pm c\delta u^2=p$ is soluble.”

\unitid{dirichlet.v1.p04b.en.par.014}However that may be, one may remark that the theorem which we are about to state is rigorously proved, by what precedes, for all values of $b$ such that the formula $t^2-bu^2$ has only the single quadratic divisor $\pm(t^2-bu^2)$. Casting one's eyes over the first of the tables added to the \emph{Théorie des Nombres}, one sees that all the prime numbers $4n+3$ smaller than 136 (the limit of the table) are in this case, excepting the single number 79. It would be easy, by following the path just indicated, to prove the truth of the theorem for this value or for any other particular value of $b$, whatever $p$ may be; but in order to establish it in all its generality, it is necessary first to prove, as we have already said, that one can always satisfy the condition which we stated only a moment ago.

\begin{center}
\unitid{dirichlet.v1.p04b.en.thm.001}\textsc{Theorem I.}
\end{center}

\unitid{dirichlet.v1.p04b.en.par.015}“$b$ denoting a prime number $4n+3$, and $p$ a prime number $4n+1$ such that $\leg{b}{p}=1$, if one puts $p=\varphi^2+\psi^2$ (where $\psi$ is supposed even), one will have the following rule for deciding whether $-b$ is or is not a biquadratic residue with respect to $p$: if $\varphi$ is divisible by $b$, $-b$ will or will not be a biquadratic residue according as $b$ is of the form $8n+7$ or of this one: $8n+3$. If $\varphi$ is not divisible by $b$, one will seek a number $\chi$ such that $\chi^2\equiv p\pmod b$. This being done, $-b$ will or will not be a biquadratic residue with respect to $p$, according as one has $\leg{\chi(\chi+\psi)}{b}=1$ or $\leg{\chi(\chi+\psi)}{b}=-1$.”

\unitid{dirichlet.v1.p04b.en.par.016}By successively taking $b=3$, $b=7$, etc., one will obtain the following particular theorems, which are analogous to the second of Mr. \textsc{Gauss}'s theorems and which must be regarded as rigorously proved by the considerations that we have set out in this paragraph and in the preceding one.

\unitid{dirichlet.v1.p04b.en.par.017}“$p$ denoting a prime number $12n+1$, if one puts $p=\varphi^2+\psi^2$ (where $\psi$ is supposed even), one of the numbers $\varphi$ and $\psi$ will necessarily be divisible by 3. This being done, I say that $-3$ will or will not be a biquadratic residue with respect to $p$, according as it is $\psi$ or $\varphi$ which is divisible by 3.”

\unitid{dirichlet.v1.p04b.en.par.018}“$p$ denoting a prime number of one of these forms: $28n+1$, $28n+9$, $28n+25$, if one puts $p=\varphi^2+\psi^2$, I say that $-7$ will be a biquadratic residue with respect to $p$ if one of the numbers $\varphi$ and $\psi$ is divisible by 7, and that $-7$ will be a biquadratic non-residue with respect to $p$ if neither of these numbers is divisible by 7.”

\unitid{dirichlet.v1.p04b.en.par.019}etc.

\begin{center}
\unitid{dirichlet.v1.p04b.en.sec.005}5.
\end{center}

\unitid{dirichlet.v1.p04b.en.par.020}Let us denote by $a$ a prime number $4n+1$, and by $p$ a prime number likewise $4n+1$, and moreover capable of being put under the form $t^2-au^2$. We can therefore put $p=t^2-au^2$, where the number $t$, as is easy to see, is necessarily odd. One will find, as in paragraph 3, that if one denotes by $2^\nu$ the highest power of 2 that divides $u$, one has $\leg{u}{p}=\leg{2^\nu}{p}$, a relation which one will then change, as in the cited place, into this one: $\leg{u}{p}=(-1)^{\frac{p-1}{4}}$.

\unitid{dirichlet.v1.p04b.en.par.021}Let us now decompose the odd number $t$ into its simple factors, putting $t=gg'g''\ldots{}\times hh'h''\ldots$, where the prime numbers $g$, $g'$, $g''$, etc., $h$, $h'$, $h''$, etc. are such that:
\unitid{dirichlet.v1.p04b.en.eq.014}
\[
(\varepsilon)\qquad
\left\{
\begin{array}{llll}
\leg{-a}{g}=1,& \leg{-a}{g'}=1,& \leg{-a}{g''}=1,& \text{etc.}\\[0.3em]
\leg{-a}{h}=-1,& \leg{-a}{h'}=-1,& \leg{-a}{h''}=-1,& \text{etc.}
\end{array}
\right.
\]

\unitid{dirichlet.v1.p04b.en.par.022}The equation $t^2-au^2=p$ gives immediately:
\unitid{dirichlet.v1.p04b.en.eq.015}
\[
\begin{array}{llll}
\leg{-ap}{g}=1,& \leg{-ap}{g'}=1,& \leg{-ap}{g''}=1,& \text{etc.}\\[0.3em]
\leg{-ap}{h}=1,& \leg{-ap}{h'}=1,& \leg{-ap}{h''}=1,& \text{etc.}
\end{array}
\]
Comparing these relations with the preceding ones, one obtains:
\unitid{dirichlet.v1.p04b.en.eq.016}
\[
\begin{array}{llll}
\leg{p}{g}=1,& \leg{p}{g'}=1,& \leg{p}{g''}=1,& \text{etc.}\\[0.3em]
\leg{p}{h}=-1,& \leg{p}{h'}=-1,& \leg{p}{h''}=-1,& \text{etc.}
\end{array}
\]
The application of the law of reciprocity will then give, $p$ being of the form $4n+1$:
\unitid{dirichlet.v1.p04b.en.eq.017}
\[
\begin{array}{llll}
\leg{g}{p}=1,& \leg{g'}{p}=1,& \leg{g''}{p}=1,& \text{etc.}\\[0.3em]
\leg{h}{p}=-1,& \leg{h'}{p}=-1,& \leg{h''}{p}=-1,& \text{etc.}
\end{array}
\]
whence one concludes, by multiplying:
\unitid{dirichlet.v1.p04b.en.eq.018}
\[
\leg{gg'g''\ldots{}\times hh'h''\ldots}{p}=\leg{t}{p}=\pm1,
\]
the upper or lower sign taking place according as the prime numbers $h$, $h'$, $h''$, \ldots{} are even or odd in number. One can state this result in a slightly different way, by saying that $\leg{t}{p}$ is equal to the product of the second members of the relations $(\varepsilon)$.

\unitid{dirichlet.v1.p04b.en.par.023}In the relations $(\varepsilon)$, one may everywhere change $-a$ into $a$, provided that at the same time one changes the sign of the second members of those of these relations in which there occurs a prime number $g$, $g'$, $g''$, \ldots{}, $h$, $h'$, $h''$, \ldots{} of the form $4n+3$. One will thus have:
\unitid{dirichlet.v1.p04b.en.eq.019}
\[
(\zeta)\qquad
\left\{
\begin{array}{llll}
\leg{a}{g}=\pm1,& \leg{a}{g'}=\pm1,& \leg{a}{g''}=\pm1,& \text{etc.}\\[0.3em]
\leg{a}{h}=\pm1,& \leg{a}{h'}=\pm1,& \leg{a}{h''}=\pm1,& \text{etc.}
\end{array}
\right.
\]

\unitid{dirichlet.v1.p04b.en.par.024}Comparing the second member of each of these relations with the second member of the corresponding relation of table $(\varepsilon)$, one will find as many changes of sign as there are prime numbers $4n+3$ among the numbers $g$, $g'$, $g''$, \ldots{}, $h$, $h'$, $h''$, \ldots{}.

\unitid{dirichlet.v1.p04b.en.par.025}The number of changes will therefore be even when among the numbers $g$, $g'$, $g''$, \ldots{}, $h$, $h'$, $h''$, \ldots{} there is an even number of the form $4n+3$. In this case, which will occur every time that $t=gg'g''\ldots{}\times hh'h''\ldots$ is of the form $4n+1$, the product of the second members of the relations $(\zeta)$ will therefore be the same as the product of the second members of the relations $(\varepsilon)$.

\unitid{dirichlet.v1.p04b.en.par.026}On the other hand, one may, by virtue of the law of reciprocity, reverse the first members of the relations $(\zeta)$ without its being necessary to change the sign of any of the second members of these relations. The product of all the reversed expressions, that is, $\leg{gg'g''\ldots{}\times hh'h''\ldots}{a}$ or $\leg{t}{a}$, is therefore equal to the product of the second members of the relations $(\zeta)$ and consequently also, according to what has previously been seen, equal to the product of the second members of the relations $(\varepsilon)$. Now, having proved above that this last product is equal to $\leg{t}{p}$, one has:
\unitid{dirichlet.v1.p04b.en.eq.020}
\[
\leg{t}{p}=\leg{t}{a}.
\]

\unitid{dirichlet.v1.p04b.en.par.027}This result concerns the case in which $t$ is of the form $4n+1$; if one examines in a similar way the case of $t=4n+3$, one will find that one then has:
\unitid{dirichlet.v1.p04b.en.eq.021}
\[
\leg{t}{p}=-\leg{t}{a}.
\]

\end{document}
